% GENERATED FILE. Edit canonical lessons, then run npm run generate:books.
% canonical-source-hash: fnv1a64:c2b0faf5bae34709
% canonical-lessons: MR-C332-vir, MR-C332-parakram, MR-C332-saurya, MR-C332-balidan, MR-C332-saran-jane

\chapter{Courage and Sacrifice}
\label{ch:mr-courage-and-sacrifice}
\hlchaptermodality{332}

\begin{chapteropening}
\textbf{By the end of this chapter:} \emph{I can talk about heroes, their deeds and surrender.}

Complete the last lesson of chapter 332: I can talk about heroes, their deeds and surrender.
\end{chapteropening}

\section[\texorpdfstring{\mr{वीर} (vīr)}{वीर (vīr)}]{\mr{वीर} (vīr) — a hero; a warrior}

\label{lesson:MR-C332-vir}



\begin{quote}
Before the new one: say the Marathi for a rampart, then the Marathi for a bastion.
\end{quote}

\subsection*{You'll want to know: \mr{वीर}}
\textbf{\mr{वीर}} — \emph{vīr} — \textquotedblleft{}a hero; a warrior\textquotedblright{}.

\begin{quote}
\mr{गावात} \mr{त्या} \mr{वीराचं} \mr{स्मारक} \mr{आहे}.
\end{quote}

\begin{quote}
\emph{gāvāt tyā vīrācaṁ smārak āhe.}
\end{quote}

\begin{quote}
There is a memorial to that hero in the village.
\end{quote}

Masculine.

As an adjective, brave: \mr{वीर} \mr{सैनिक}.

\begin{cousinweb}
\textbf{\mr{वीर}} (\emph{vīr}) is a tatsama, taken from Sanskrit \emph{vīra}, 'a man, a hero' (Monier-Williams).
\end{cousinweb}

\begin{grammarlens}[title={What you've built}]
One more word for telling what happened.
\end{grammarlens}

\subsection*{Your turn}
\begin{itemize}
\raggedright
  \item \emph{Say it:} \emph{vīr}
  \item \emph{Say it:} \emph{vīr}, once more
  \item \emph{Say it:} \emph{gāvāt tyā vīrācaṁ smārak āhe.} — There is a memorial to that hero in the village.
  \item \emph{Recall:} say the Marathi for a rampart, then the Marathi for a bastion, then say \emph{vīr} again
\end{itemize}

\begin{tcolorbox}[breakable,colback=teal!4,colframe=teal!35!black,title={Before you move on}]
What does \mr{वीर} mean? (\textquotedblleft{}A hero; a warrior\textquotedblright{}.) Say it once more.
\end{tcolorbox}
\section[\texorpdfstring{\mr{पराक्रम} (parākram)}{पराक्रम (parākram)}]{\mr{पराक्रम} (parākram) — a heroic deed; valour}

\label{lesson:MR-C332-parakram}



\begin{quote}
Before the new one: say the Marathi for a bastion, then the Marathi for a hero.
\end{quote}

\subsection*{You'll want to know: \mr{पराक्रम}}
\textbf{\mr{पराक्रम}} — \emph{parākram} — \textquotedblleft{}a heroic deed; valour\textquotedblright{}.

\begin{quote}
\mr{लोक} \mr{आजही} \mr{त्याच्या} \mr{पराक्रमाच्या} \mr{कथा} \mr{सांगतात}.
\end{quote}

\begin{quote}
\emph{lok ājhī tyācyā parākramācyā kathā sāṅgtāt.}
\end{quote}

\begin{quote}
Stories of his valour are still told today.
\end{quote}

Masculine.

\mr{पराक्रम} \mr{गाजवणे}, to perform great deeds.

\begin{cousinweb}
\textbf{\mr{पराक्रम}} (\emph{parākram}) is a tatsama, taken from Sanskrit \emph{parākrama}, 'a bold advance, heroism' (Monier-Williams).
\end{cousinweb}

\begin{grammarlens}[title={What you've built}]
One more word for telling what happened.
\end{grammarlens}

\subsection*{Your turn}
\begin{itemize}
\raggedright
  \item \emph{Say it:} \emph{parākram}
  \item \emph{Say it:} \emph{parākram}, once more
  \item \emph{Say it:} \emph{lok ājhī tyācyā parākramācyā kathā sāṅgtāt.} — Stories of his valour are still told today.
  \item \emph{Recall:} say the Marathi for a bastion, then the Marathi for a hero, then say \emph{parākram} again
\end{itemize}

\begin{tcolorbox}[breakable,colback=teal!4,colframe=teal!35!black,title={Before you move on}]
What does \mr{पराक्रम} mean? (\textquotedblleft{}A heroic deed; valour\textquotedblright{}.) Say it once more.
\end{tcolorbox}
\section[\texorpdfstring{\mr{शौर्य} (śaurya)}{शौर्य (śaurya)}]{\mr{शौर्य} (śaurya) — bravery}

\label{lesson:MR-C332-saurya}



\begin{quote}
Before the new one: say the Marathi for a hero, then the Marathi for a heroic deed.
\end{quote}

\subsection*{You'll want to know: \mr{शौर्य}}
\textbf{\mr{शौर्य}} — \emph{śaurya} — \textquotedblleft{}bravery\textquotedblright{}.

\begin{quote}
\mr{सैनिकांच्या} \mr{शौर्यामुळे} \mr{गाव} \mr{वाचलं}.
\end{quote}

\begin{quote}
\emph{sainikāñcyā śauryāmuḷe gāv vāclaṁ.}
\end{quote}

\begin{quote}
The village was saved by the soldiers' bravery.
\end{quote}

Neuter.

\mr{शौर्य} \mr{दाखवणे}, to show bravery.

\begin{cousinweb}
\textbf{\mr{शौर्य}} (\emph{śaurya}) is a tatsama, taken from Sanskrit \emph{śaurya}, 'heroism, valour' (Monier-Williams), built on \emph{śūra}, 'a hero'.
\end{cousinweb}

\begin{grammarlens}[title={What you've built}]
One more word for telling what happened.
\end{grammarlens}

\subsection*{Your turn}
\begin{itemize}
\raggedright
  \item \emph{Say it:} \emph{śaurya}
  \item \emph{Say it:} \emph{śaurya}, once more
  \item \emph{Say it:} \emph{sainikāñcyā śauryāmuḷe gāv vāclaṁ.} — The village was saved by the soldiers' bravery.
  \item \emph{Recall:} say the Marathi for a hero, then the Marathi for a heroic deed, then say \emph{śaurya} again
\end{itemize}

\begin{tcolorbox}[breakable,colback=teal!4,colframe=teal!35!black,title={Before you move on}]
What does \mr{शौर्य} mean? (\textquotedblleft{}Bravery\textquotedblright{}.) Say it once more.
\end{tcolorbox}
\section[\texorpdfstring{\mr{बलिदान} (balidān)}{बलिदान (balidān)}]{\mr{बलिदान} (balidān) — a sacrifice}

\label{lesson:MR-C332-balidan}



\begin{quote}
Before the new one: say the Marathi for a heroic deed, then the Marathi for bravery.
\end{quote}

\subsection*{You'll want to know: \mr{बलिदान}}
\textbf{\mr{बलिदान}} — \emph{balidān} — \textquotedblleft{}a sacrifice\textquotedblright{}.

\begin{quote}
\mr{देशासाठी} \mr{अनेकांनी} \mr{बलिदान} \mr{दिलं}.
\end{quote}

\begin{quote}
\emph{deśāsāṭhī anekānnī balidān dilaṁ.}
\end{quote}

\begin{quote}
Many gave their lives for the country.
\end{quote}

Neuter.

\mr{बलिदान} \mr{देणे}, to give one's life for a cause.

\begin{cousinweb}
\textbf{\mr{बलिदान}} (\emph{balidān}) is a tatsama, taken from Sanskrit \emph{balidāna}, 'the giving of an offering' (Monier-Williams).
\end{cousinweb}

\begin{grammarlens}[title={What you've built}]
One more word for telling what happened.
\end{grammarlens}

\subsection*{Your turn}
\begin{itemize}
\raggedright
  \item \emph{Say it:} \emph{balidān}
  \item \emph{Say it:} \emph{balidān}, once more
  \item \emph{Say it:} \emph{deśāsāṭhī anekānnī balidān dilaṁ.} — Many gave their lives for the country.
  \item \emph{Recall:} say the Marathi for a heroic deed, then the Marathi for bravery, then say \emph{balidān} again
\end{itemize}

\begin{tcolorbox}[breakable,colback=teal!4,colframe=teal!35!black,title={Before you move on}]
What does \mr{बलिदान} mean? (\textquotedblleft{}A sacrifice\textquotedblright{}.) Say it once more.
\end{tcolorbox}
\section[\texorpdfstring{\mr{शरण} \mr{जाणे} (śaraṇ jāṇe)}{शरण जाणे (śaraṇ jāṇe)}]{\mr{शरण} \mr{जाणे} (śaraṇ jāṇe) — to surrender}

\label{lesson:MR-C332-saran-jane}



\begin{quote}
Before the new one: say the Marathi for bravery, then the Marathi for a sacrifice.
\end{quote}

\subsection*{You'll want to know: \mr{शरण} \mr{जाणे}}
\textbf{\mr{शरण} \mr{जाणे}} — \emph{śaraṇ jāṇe} — \textquotedblleft{}to surrender\textquotedblright{}.

\begin{quote}
\mr{शेवटी} \mr{शत्रू} \mr{राजाला} \mr{शरण} \mr{गेला}.
\end{quote}

\begin{quote}
\emph{śevṭī śatrū rājālā śaraṇ gelā.}
\end{quote}

\begin{quote}
In the end the enemy surrendered to the king.
\end{quote}

Past (he …): \mr{तो} \mr{शरण} \mr{गेला}, \emph{to śaraṇ gelā}.

The one surrendered to takes -\mr{ला}: \mr{राजाला} \mr{शरण} \mr{जाणे}.

\begin{cousinweb}
\textbf{\mr{शरण} \mr{जाणे}} (\emph{śaraṇ jāṇe}) joins \emph{śaraṇ}, a tatsama taken from Sanskrit \emph{śaraṇa}, 'shelter, refuge' (Monier-Williams), to \emph{jāṇe}, 'to go': to go for refuge.
\end{cousinweb}

\begin{grammarlens}[title={What you've built}]
One more word for telling what happened.
\end{grammarlens}

\subsection*{Your turn}
\begin{itemize}
\raggedright
  \item \emph{Say it:} \emph{śaraṇ jāṇe}
  \item \emph{Say it:} \emph{śaraṇ jāṇe}, once more
  \item \emph{Say it:} \emph{śevṭī śatrū rājālā śaraṇ gelā.} — In the end the enemy surrendered to the king.
  \item \emph{Recall:} say the Marathi for bravery, then the Marathi for a sacrifice, then say \emph{śaraṇ jāṇe} again
\end{itemize}

\begin{tcolorbox}[breakable,colback=teal!4,colframe=teal!35!black,title={Before you move on}]
What does \mr{शरण} \mr{जाणे} mean? (\textquotedblleft{}To surrender\textquotedblright{}.) Say it once more.
\end{tcolorbox}
